%=====================================================================
\section{Zweck, Umfang und Lesehinweise}\label{sec:scope}
%=====================================================================

Dieses Dokument ist die verbindliche technische Spezifikation für die Implementierung von
\textbf{FourierGrating}, einem Löser für die Beugung elektromagnetischer Wellen an periodischen
Strukturen nach der Fourier-Modal-Methode (FMM). Es richtet sich an einen KI-Agenten (oder
Entwickler), der die Implementierung ohne weitere Rückfragen beginnen soll. Alle Formeln sind
hergeleitet oder auf geprüfte Quellen zurückgeführt; Stellen, an denen die Vorlage
(Nevière/Popov, \emph{Light Propagation in Periodic Media}, 2003) fehlerhaft ist, sind in
Abschnitt~\ref{sec:errata} aufgeführt und im Text bereits korrigiert.

\subsection{Funktionsumfang}
\begin{center}
\begin{tabular}{lll}
\toprule
Variante & Periodizität & Material \\
\midrule
1D-isotrop    & $x$ (eine Fouriertransformation, FT) & $\varepsilon$ skalar, $\mu=1$ \\
1D-anisotrop  & $x$                                  & $\overline{\overline\varepsilon}$ voller komplexer $3\times3$-Tensor \\
2D-isotrop    & $x$ und $z$ (zwei FT)               & $\varepsilon$ skalar \\
2D-anisotrop  & $x$ und $z$                          & $\overline{\overline\varepsilon}$ voller Tensor \\
\bottomrule
\end{tabular}
\end{center}
Für jede Variante gibt es zwei Löser für den modulierten Bereich:
\begin{enumerate}[nosep]
  \item \textbf{RK}: Differentialmethode, numerische Integration der Propagationsgleichung
        $dF/dy=M(y)F$ mit einem \emph{impliziten} Runge--Kutta-Verfahren (Gauß--Legendre, optional
        Radau IIA), stabilisiert durch die S-Matrix.
  \item \textbf{RCWA}: Zerlegung in Scheiben mit $y$-unabhängigem $M$ (Treppenapproximation),
        Eigenwertzerlegung je Scheibe, stabile Scheiben-S-Matrix.
\end{enumerate}
Beide Löser verwenden dieselbe Fourierfaktorisierung (FFF, \enquote{Fast Fourier Factorization}
mit Normalenvektorfeldern), dieselben Randbedingungen, dieselbe S-Matrix-Verkettung und dieselbe
Nahfeldrekonstruktion. Das Nahfeld wird jeweils in einem \textbf{horizontalen} ($y=\text{const}$)
und einem \textbf{vertikalen} Schnitt ($z=\text{const}$, Ebene $x$--$y$) ausgegeben.

Zielsystem: AlmaLinux 10 (x86\_64, GCC 14), C++20, CUDA (NVIDIA RTX 6000 Ada,
Compute Capability 8.9, 48\,GB). Ein reiner CPU-Pfad ist Pflicht (Referenz, Tests, CI).

\subsection{Nicht-Ziele (vorerst)}
Magnetische Anisotropie ($\overline{\overline\mu}\neq\mathbb 1$) wird in den Formeln mitgeführt, aber erst
nach Meilenstein M9 implementiert. Nichtlineare Optik, Photonische-Kristall-Bandstrukturen,
Zylinderkoordinaten und Polologie (Buchkapitel X--XV) sind nicht Teil dieses Plans.

\subsection{Lesehinweise für den Agenten}
\begin{itemize}[nosep]
  \item Abschnitte \ref{sec:conv}--\ref{sec:field} enthalten die Mathematik, Abschnitt
        \ref{sec:arch} die Software-Architektur, Abschnitt \ref{sec:validation} die Tests mit
        Referenzwerten, Abschnitt \ref{sec:milestones} die Reihenfolge der Arbeit.
  \item Jede Formel, die implementiert wird, bekommt einen Unit-Test (Abschnitt~\ref{sec:validation}).
  \item Bei Widerspruch zwischen diesem Dokument und dem Buch gilt dieses Dokument.
  \item Gleichungsnummern des Buches werden als \enquote{Buch (IV.21)} zitiert.
\end{itemize}

%=====================================================================
\section{Konventionen und normierte Einheiten}\label{sec:conv}
%=====================================================================

\subsection{Zeitabhängigkeit, Koordinaten, Schichtung}
\begin{itemize}[nosep]
  \item Zeitabhängigkeit $e^{-i\omega t}$; Felder sind komplexe Amplituden.
  \item Periodizität in $x$ (Periode $d_x$) und -- im 2D-Fall -- in $z$ (Periode $d_z$). Die
        Schichtung (Stapelrichtung) ist $y$. Bei 1D-Gittern ist die Struktur in $z$ invariant.
  \item Das \textbf{Substrat} (Medium 1) belegt $y<0$, das \textbf{Superstrat} (Medium $M$) belegt
        $y>y_{\text{top}}$. Dazwischen liegen Schichten $j=2,\dots,M-1$ mit Grenzen
        $0=y_1<y_2<\dots<y_{M-1}=y_{\text{top}}$. Licht fällt von oben ein.
  \item \enquote{Abwärts} (down, Amplitude $a$) heißt Ausbreitung bzw.\ Abklingen in $-y$,
        \enquote{aufwärts} (up, Amplitude $b$) in $+y$.
\end{itemize}

\subsection{Normierung}
Mit $k_0=\omega/c=2\pi/\lambda_0$ und $Z_0=\sqrt{\mu_0/\varepsilon_0}$ werden eingeführt:
\begin{equation}
  \tilde{\mathbf H}=Z_0\mathbf H,\qquad
  \hat x=k_0x,\ \hat y=k_0y,\ \hat z=k_0z,\qquad
  \hat\alpha=\alpha/k_0,\ \hat\beta=\beta/k_0,\ \hat\gamma=\gamma/k_0 .
  \label{eq:norm}
\end{equation}
Mit relativen Materialtensoren $\overline{\overline\varepsilon}_r$, $\overline{\overline\mu}_r$ lauten die
Maxwellgleichungen dann
\begin{equation}
  \widehat{\operatorname{rot}}\,\mathbf E=i\,\overline{\overline\mu}_r\tilde{\mathbf H},\qquad
  \widehat{\operatorname{rot}}\,\tilde{\mathbf H}=-i\,\overline{\overline\varepsilon}_r\mathbf E .
  \label{eq:maxwell_norm}
\end{equation}
Herleitung: $\operatorname{rot}\mathbf E=i\omega\mu_0\mu_r\mathbf H=i\omega\sqrt{\mu_0\varepsilon_0}\,\mu_r\tilde{\mathbf H}
=ik_0\mu_r\tilde{\mathbf H}$ und
$\operatorname{rot}\tilde{\mathbf H}=Z_0(-i\omega\varepsilon_0\varepsilon_r\mathbf E)=-ik_0\varepsilon_r\mathbf E$;
anschließend $\partial/\partial x=k_0\,\partial/\partial\hat x$.
\textbf{Im Code werden ausschließlich normierte Größen verwendet} (Hut wird im Folgenden weggelassen;
alle Längen in Einheiten von $1/k_0$). Benutzereingaben (µm, Grad) werden an der Schnittstelle
umgerechnet. Damit entsprechen alle Buchformeln dem Fall $\omega=1$, $\mu_0=\varepsilon_0=1$.

\subsection{Wellenzahlen, Wurzelzweig}
\begin{equation}
  \alpha_n=\alpha_0+nK_x,\quad K_x=\frac{2\pi}{k_0d_x},\qquad
  \gamma_m=\gamma_0+mK_z,\quad K_z=\frac{2\pi}{k_0d_z}
  \label{eq:alpha_gamma}
\end{equation}
(im 1D-Fall $\gamma_m\equiv\gamma_0$). In einem isotropen homogenen Medium mit $\varepsilon$:
\begin{equation}
  \beta_{nm}=\sqrt{\varepsilon-\alpha_n^2-\gamma_m^2},\qquad
  \operatorname{Im}\beta\ge0,\ \text{und}\ \operatorname{Re}\beta\ge0\ \text{falls}\ \operatorname{Im}\beta=0 .
  \label{eq:beta_branch}
\end{equation}
\textbf{Implementierung:}\\ \verb!b = sqrt(z); if (b.imag()<0 || (b.imag()==0 && b.real()<0)) b = -b;!
\textbf{Rayleigh-Anomalie:} Falls $|\beta|<\eta_{\text{RA}}$ (Standard $10^{-10}$), wird
$\beta\to\beta+i\eta_{\text{RA}}$ gesetzt und eine Warnung protokolliert.

\subsection{Einfallende Welle}
$\Phi$ = Winkel zwischen $\mathbf k$ und der $y$-Achse, $\Psi$ = Azimut der Projektion auf die
$xz$-Ebene gegen $x$, $\delta$ = Polarisationswinkel (Buch Kap.~IX, korrigiert):
\begin{equation}
  \alpha_0=\nu_M\sin\Phi\cos\Psi,\qquad \beta_0=\nu_M\cos\Phi,\qquad \gamma_0=\nu_M\sin\Phi\sin\Psi,
  \qquad \mathbf k^{(i)}=(\alpha_0,-\beta_0,\gamma_0),
  \label{eq:incidence}
\end{equation}
($\nu_M$: reelle Brechzahl des Superstrats; Buch (IX.1) hat fälschlich $\cos\Psi$ in $\beta_0$),
\begin{equation}
  \mathbf A=\begin{pmatrix}\sin\Psi\cos\delta+\sin\delta\cos\Phi\cos\Psi\\ \sin\delta\sin\Phi\\
  \cos\Phi\sin\Psi\sin\delta-\cos\delta\cos\Psi\end{pmatrix},\qquad |\mathbf A|=1,\ \mathbf A\cdot\mathbf k^{(i)}=0 .
  \label{eq:polarization}
\end{equation}
$\delta=0$: $\mathbf A\perp$ Einfallsebene (\enquote{s}); $\delta=90^\circ$: in der Einfallsebene (\enquote{p}).
Bei Normaleinfall gilt $\Psi=\pi$, dann ist $\delta$ der Winkel zwischen $\mathbf A$ und $z$.
1D-in-plane: $\Psi=0$ bzw.\ $\pi$, $\gamma_0=0$; TE $\equiv\delta=0$ ($\mathbf E\parallel z$), TM $\equiv\delta=90^\circ$.

%=====================================================================
\section{Grundgleichungen im Fourierraum}\label{sec:equations}
%=====================================================================

\subsection{Komponentenform}
Aus \eqref{eq:maxwell_norm} (Buch (V.14)--(V.19) mit $\omega=1$):
\begin{align}
  \partial_yE_z-\partial_zE_y&=iB_x, & \partial_yH_z-\partial_zH_y&=-iD_x,\label{eq:mx1}\\
  \partial_zE_x-\partial_xE_z&=iB_y, & \partial_zH_x-\partial_xH_z&=-iD_y,\label{eq:mx2}\\
  \partial_xE_y-\partial_yE_x&=iB_z, & \partial_xH_y-\partial_yH_x&=-iD_z,\label{eq:mx3}
\end{align}
mit $\mathbf D=\overline{\overline\varepsilon}_r\mathbf E$, $\mathbf B=\overline{\overline\mu}_r\mathbf H$ (hier stets
$\mathbf H\equiv\tilde{\mathbf H}$).

\subsection{Pseudoperiodische Entwicklung}
Wegen der Pseudoperiodizität (Buch I.2.3) gilt für jede Feldkomponente
\begin{equation}
  u(x,y,z)=\sum_{(n,m)\in\mathcal N}u_{nm}(y)\,e^{i(\alpha_nx+\gamma_mz)},
  \label{eq:pseudo}
\end{equation}
mit der Abschneidemenge $\mathcal N$ (Abschnitt~\ref{sec:fourier}). Mit $\alpha=\operatorname{diag}(\alpha_n)$,
$\gamma=\operatorname{diag}(\gamma_m)$ (im Einzelindexraum, Abschnitt~\ref{sec:indexmap}) wird
$\partial_x\to i\alpha$, $\partial_z\to i\gamma$. Produkte $\varepsilon E$ werden durch die
Faktorisierung aus Abschnitt~\ref{sec:fff} ersetzt:
\begin{equation}
  [\mathbf D]=Q_\varepsilon[\mathbf E],\qquad [\mathbf B]=Q_\mu[\mathbf H],
  \label{eq:Qdef}
\end{equation}
mit $3\times3$-Blockmatrizen $Q_\varepsilon=(Q_{\varepsilon,ij})_{i,j\in\{x,y,z\}}$, jeder Block
$n_f\times n_f$ ($n_f=|\mathcal N|$). Ohne magnetische Anisotropie ist $Q_\mu=\mathbb 1$.

\subsection{Elimination von $E_y$, $H_y$}
Aus \eqref{eq:mx2} (rechte Spalte): $i\gamma[H_x]-i\alpha[H_z]=-i[D_y]$, also
$[D_y]=\alpha[H_z]-\gamma[H_x]$, und mit \eqref{eq:Qdef}
\begin{equation}
  [E_y]=Q_{\varepsilon,yy}^{-1}\big(\alpha[H_z]-\gamma[H_x]-Q_{\varepsilon,yx}[E_x]-Q_{\varepsilon,yz}[E_z]\big),
  \label{eq:Ey}
\end{equation}
und analog aus \eqref{eq:mx2} (links): $i[B_y]=i\gamma[E_x]-i\alpha[E_z]$:
\begin{equation}
  [H_y]=Q_{\mu,yy}^{-1}\big(\gamma[E_x]-\alpha[E_z]-Q_{\mu,yx}[H_x]-Q_{\mu,yz}[H_z]\big).
  \label{eq:Hy}
\end{equation}

\subsection{Propagationsgleichung}
Mit dem Zustandsvektor $F=([E_x],[E_z],[H_x],[H_z])^T\in\mathbb C^{4n_f}$ gilt
\begin{equation}
  \frac{dF}{dy}=M(y)\,F .
  \label{eq:ode}
\end{equation}
Herleitung der Zeilen: $\partial_yE_x=i\alpha E_y-iB_z$ \eqref{eq:mx3}, $\partial_yE_z=i\gamma E_y+iB_x$
\eqref{eq:mx1}, $\partial_yH_x=i\alpha H_y+iD_z$ \eqref{eq:mx3}, $\partial_yH_z=i\gamma H_y-iD_x$
\eqref{eq:mx1}; Einsetzen von \eqref{eq:Ey}, \eqref{eq:Hy} und
$B_x=Q_{\mu,xx}H_x+Q_{\mu,xy}H_y+Q_{\mu,xz}H_z$ usw. ergibt
(mit $\mathcal E\equiv Q_{\varepsilon,yy}^{-1}$, $\mathcal U\equiv Q_{\mu,yy}^{-1}$, Buch (V.22') mit $\omega=1$):
\begin{equation}
\begin{aligned}
  M_{11}&=-i\alpha\mathcal EQ_{\varepsilon,yx}-iQ_{\mu,zy}\mathcal U\gamma, &
  M_{12}&=-i\alpha\mathcal EQ_{\varepsilon,yz}+iQ_{\mu,zy}\mathcal U\alpha,\\
  M_{13}&=iQ_{\mu,zy}\mathcal UQ_{\mu,yx}-i\alpha\mathcal E\gamma-iQ_{\mu,zx}, &
  M_{14}&=iQ_{\mu,zy}\mathcal UQ_{\mu,yz}+i\alpha\mathcal E\alpha-iQ_{\mu,zz},\\
  M_{21}&=-i\gamma\mathcal EQ_{\varepsilon,yx}+iQ_{\mu,xy}\mathcal U\gamma, &
  M_{22}&=-i\gamma\mathcal EQ_{\varepsilon,yz}-iQ_{\mu,xy}\mathcal U\alpha,\\
  M_{23}&=-iQ_{\mu,xy}\mathcal UQ_{\mu,yx}-i\gamma\mathcal E\gamma+iQ_{\mu,xx}, &
  M_{24}&=-iQ_{\mu,xy}\mathcal UQ_{\mu,yz}+i\gamma\mathcal E\alpha+iQ_{\mu,xz},\\
  M_{31}&=-iQ_{\varepsilon,zy}\mathcal EQ_{\varepsilon,yx}+i\alpha\mathcal U\gamma+iQ_{\varepsilon,zx}, &
  M_{32}&=-iQ_{\varepsilon,zy}\mathcal EQ_{\varepsilon,yz}-i\alpha\mathcal U\alpha+iQ_{\varepsilon,zz},\\
  M_{33}&=-i\alpha\mathcal UQ_{\mu,yx}-iQ_{\varepsilon,zy}\mathcal E\gamma, &
  M_{34}&=-i\alpha\mathcal UQ_{\mu,yz}+iQ_{\varepsilon,zy}\mathcal E\alpha,\\
  M_{41}&=iQ_{\varepsilon,xy}\mathcal EQ_{\varepsilon,yx}+i\gamma\mathcal U\gamma-iQ_{\varepsilon,xx}, &
  M_{42}&=iQ_{\varepsilon,xy}\mathcal EQ_{\varepsilon,yz}-i\gamma\mathcal U\alpha-iQ_{\varepsilon,xz},\\
  M_{43}&=-i\gamma\mathcal UQ_{\mu,yx}+iQ_{\varepsilon,xy}\mathcal E\gamma, &
  M_{44}&=-i\gamma\mathcal UQ_{\mu,yz}-iQ_{\varepsilon,xy}\mathcal E\alpha .
\end{aligned}
\label{eq:Mgeneral}
\end{equation}
Für $Q_\mu=\mathbb 1$ ($\mathcal U=\mathbb 1$, $Q_{\mu,ij}=0$ für $i\neq j$) vereinfacht sich das zu
\begin{equation}
\begin{aligned}
  M_{11}&=-i\alpha\mathcal EQ_{\varepsilon,yx}, & M_{12}&=-i\alpha\mathcal EQ_{\varepsilon,yz},
  & M_{13}&=-i\alpha\mathcal E\gamma, & M_{14}&=i\alpha\mathcal E\alpha-i\mathbb 1,\\
  M_{21}&=-i\gamma\mathcal EQ_{\varepsilon,yx}, & M_{22}&=-i\gamma\mathcal EQ_{\varepsilon,yz},
  & M_{23}&=-i\gamma\mathcal E\gamma+i\mathbb 1, & M_{24}&=i\gamma\mathcal E\alpha,
\end{aligned}
\label{eq:Mmu1a}
\end{equation}
\begin{equation}
\begin{aligned}
  M_{31}&=i\big(Q_{\varepsilon,zx}-Q_{\varepsilon,zy}\mathcal EQ_{\varepsilon,yx}+\alpha\gamma\big), &
  M_{32}&=i\big(Q_{\varepsilon,zz}-Q_{\varepsilon,zy}\mathcal EQ_{\varepsilon,yz}-\alpha^2\big),\\
  M_{33}&=-iQ_{\varepsilon,zy}\mathcal E\gamma, & M_{34}&=iQ_{\varepsilon,zy}\mathcal E\alpha,\\
  M_{41}&=i\big(Q_{\varepsilon,xy}\mathcal EQ_{\varepsilon,yx}-Q_{\varepsilon,xx}+\gamma^2\big), &
  M_{42}&=i\big(Q_{\varepsilon,xy}\mathcal EQ_{\varepsilon,yz}-Q_{\varepsilon,xz}-\gamma\alpha\big),\\
  M_{43}&=iQ_{\varepsilon,xy}\mathcal E\gamma, & M_{44}&=-iQ_{\varepsilon,xy}\mathcal E\alpha .
\end{aligned}
\label{eq:Mmu1b}
\end{equation}
\begin{tcolorbox}[title=Prüfpunkt (Unit-Test \texttt{test\_propagation\_matrix})]
In einem homogenen isotropen Medium ($Q_{\varepsilon,ij}=\varepsilon\delta_{ij}\mathbb 1$) müssen die
analytischen Moden aus Abschnitt~\ref{sec:homog} Eigenvektoren von $M$ zu den Eigenwerten
$\lambda=\pm i\beta_{nm}$ sein. Achtung: Buch (VIII.3) lässt in $M_{32}$ den Term $-i\alpha^2$ weg
(falsch); \eqref{eq:Mmu1b} ist korrekt.
\end{tcolorbox}

\subsection{Spezialfälle}
\paragraph{1D, in-plane ($\gamma=0$), isotrop, $N_z=0$:} Dann $Q_{\varepsilon,xz}=Q_{\varepsilon,zx}=Q_{\varepsilon,yz}=
Q_{\varepsilon,zy}=0$, $Q_{\varepsilon,zz}=\tp{\varepsilon}$, und \eqref{eq:ode} zerfällt in
\begin{align}
  \text{TE: }&\frac{d}{dy}\begin{pmatrix}[E_z]\\{}[H_x]\end{pmatrix}=
   \begin{pmatrix}0&i\mathbb 1\\ i(\tp{\varepsilon}-\alpha^2)&0\end{pmatrix}\begin{pmatrix}[E_z]\\{}[H_x]\end{pmatrix},
   \label{eq:TE}\\
  \text{TM: }&\frac{d}{dy}\begin{pmatrix}[E_x]\\{}[H_z]\end{pmatrix}=
   \begin{pmatrix}-i\alpha\mathcal EQ_{\varepsilon,yx}& i\alpha\mathcal E\alpha-i\mathbb 1\\
   i(Q_{\varepsilon,xy}\mathcal EQ_{\varepsilon,yx}-Q_{\varepsilon,xx})&-iQ_{\varepsilon,xy}\mathcal E\alpha\end{pmatrix}
   \begin{pmatrix}[E_x]\\{}[H_z]\end{pmatrix}.\label{eq:TM}
\end{align}
Mit $Q_{\varepsilon,xx}=A+\Li$, $Q_{\varepsilon,yy}=\tp{\varepsilon}-A$, $Q_{\varepsilon,xy}=Q_{\varepsilon,yx}=B$
($\Delta=\tp\varepsilon-\Li$, $A=\Delta\tp{c^2}$, $B=\Delta\tp{cs}$, $c=\cos\theta$, $s=\sin\theta$,
$\tan\theta=g'$) ist \eqref{eq:TM} identisch mit Buch (IV.21), (IV.22). Diese beiden
$2n_f$-Systeme sind der erste Implementierungsschritt (Meilenstein M2/M3).
\paragraph{1D konisch:} volles $4n_f$-System mit $\gamma=\gamma_0\mathbb 1$.
\paragraph{2D:} volles $4n_f$-System mit $n_f=(2N_x+1)(2N_z+1)$.

%=====================================================================
\section{Fourier-Darstellung und Toeplitz-Matrizen}\label{sec:fourier}
%=====================================================================

\subsection{Abschneidemenge und Einzelindex}\label{sec:indexmap}
1D: $\mathcal N=\{n:\ -N_x\le n\le N_x\}$, $n_f=2N_x+1$, Index $r=n+N_x$ (0-basiert).
2D (rechteckige Abschneidung): $\mathcal N=\{(n,m):|n|\le N_x,|m|\le N_z\}$,
\begin{equation}
  r=(m+N_z)+(2N_z+1)(n+N_x),\qquad
  n=-N_x+\Big\lfloor\frac{r}{2N_z+1}\Big\rfloor,\quad m=r-(2N_z+1)(n+N_x)-N_z
  \label{eq:indexmap}
\end{equation}
(0-basiert; das Buch (IX.14) dividiert fälschlich durch $N+1$). Optional später: kreisförmige
Abschneidung $(n/N_x)^2+(m/N_z)^2\le1$; dann ist die Abbildung eine Tabelle. Der Code darf nur
über die Klasse \verb|IndexMap| auf $(n,m)\leftrightarrow r$ zugreifen.

\subsection{Fourierkoeffizienten}
Für eine $d_x$-periodische (bzw.\ $(d_x,d_z)$-periodische) Funktion $f$ gilt
\begin{equation}
  f(x,z)=\sum_{p,q}f_{pq}\,e^{i(pK_xx+qK_zz)},\qquad
  f_{pq}=\frac{1}{d_xd_z}\iint_{\text{Zelle}}f(x,z)\,e^{-i(pK_xx+qK_zz)}\,dx\,dz .
  \label{eq:fcoef}
\end{equation}
Für Toeplitz-Matrizen werden $|p|\le2N_x$, $|q|\le2N_z$ benötigt.

\paragraph{Intervall (1D).} $f=\chi_{[x_1,x_2]}$ (Periode $d$, $K=2\pi/d$):
\begin{equation}
  f_0=\frac{x_2-x_1}{d},\qquad f_p=\frac{e^{-ipKx_2}-e^{-ipKx_1}}{-ipKd}\quad(p\ne0).
  \label{eq:interval}
\end{equation}
Stückweise konstante Funktionen $=\sum_j\varepsilon_j\chi_{I_j}$ (Linearität).

\paragraph{Rechteck (2D).} Produkt zweier Intervallformeln.

\paragraph{Polygon (2D, exakt).} Für ein einfaches Polygon $P$ mit Ecken $\mathbf v_1,\dots,\mathbf v_L$
(gegen den Uhrzeigersinn), Kanten $\mathbf t_j=\mathbf v_{j+1}-\mathbf v_j$, Länge $\ell_j$, Mittelpunkt
$\mathbf c_j$, äußerer Normale $\mathbf n_j$ und $\mathbf k=(pK_x,qK_z)\neq0$:
\begin{equation}
  \iint_Pe^{-i\mathbf k\cdot\mathbf r}d^2r=\frac{i}{|\mathbf k|^2}\sum_{j=1}^{L}(\mathbf k\cdot\mathbf n_j)\,\ell_j\,
  e^{-i\mathbf k\cdot\mathbf c_j}\,\operatorname{sinc}\Big(\frac{\mathbf k\cdot\mathbf t_j}{2}\Big),\qquad
  \operatorname{sinc}u=\frac{\sin u}{u},
  \label{eq:polygon}
\end{equation}
für $\mathbf k=0$ die Fläche. Herleitung: Gaußscher Satz mit
\[
  e^{-i\mathbf k\cdot\mathbf r}=\frac{i}{|\mathbf k|^2}\nabla\cdot\big(\mathbf k\,e^{-i\mathbf k\cdot\mathbf r}\big),\qquad
  \int_0^1e^{-i\mathbf k\cdot(\mathbf v_j+s\mathbf t_j)}\ell_j\,ds=\ell_j\,e^{-i\mathbf k\cdot\mathbf c_j}\operatorname{sinc}(\mathbf k\cdot\mathbf t_j/2).
\]
Kreise/Ellipsen: analytisch über Bessel-Funktion ($J_1$) oder als feines Polygon.

\paragraph{Allgemeine Funktionen.} Abtastung auf einem feinen Gitter ($\ge 8\times$ die höchste
benötigte Ordnung) und FFT. \textbf{Nur für stetige Funktionen} (z.\,B.\ Normalenvektorkomponenten);
für unstetige Funktionen (Permittivität) möglichst analytische Formeln, sonst Aliasing-Fehler
$O(1/N_{\text{grid}})$.

\subsection{Toeplitz-Matrizen}
\begin{equation}
  \tp{f}_{rr'}=f_{n-n',\,m-m'},\qquad r\leftrightarrow(n,m),\ r'\leftrightarrow(n',m').
  \label{eq:toeplitz}
\end{equation}
1D: gewöhnliche Toeplitz-Matrix; 2D: block-Toeplitz mit Toeplitz-Blöcken. Rechenregeln
(Buch Anh.~II; im abgeschnittenen Raum nur näherungsweise):
$\tp{1}=\mathbb 1$, $\tp{f+g}=\tp f+\tp g$, $\tp{fg}\approx\tp f\tp g$ falls keine gleichzeitigen Sprünge.
\begin{tcolorbox}[title=Numerischer Befund (Kapitelprüfung)]
$\tp f\tp g-\tp{fg}$ ist im Inneren der Matrix auf Maschinengenauigkeit null, an den Rändern
(höchste Ordnungen) bleibt ein Fehler $O(1)$, der mit $N$ \emph{nicht} abnimmt. Deshalb nie
$\tp f\tp g$ durch $\tp{fg}$ ersetzen, wenn eine der Funktionen unstetig ist, außer die Spezifikation
verlangt es ausdrücklich.
\end{tcolorbox}
